% Options for packages loaded elsewhere
\PassOptionsToPackage{unicode}{hyperref}
\PassOptionsToPackage{hyphens}{url}
\documentclass[
  11pt,
  letterpaper,
]{article}
\usepackage{xcolor}
\usepackage[left=0.82in,right=0.82in,top=0.73in,bottom=0.77in]{geometry}
\usepackage{amsmath,amssymb}
\setcounter{secnumdepth}{5}
\usepackage{iftex}
\ifPDFTeX
  \usepackage[T1]{fontenc}
  \usepackage[utf8]{inputenc}
  \usepackage{textcomp} % provide euro and other symbols
\else % if luatex or xetex
  \usepackage{unicode-math} % this also loads fontspec
  \defaultfontfeatures{Scale=MatchLowercase}
  \defaultfontfeatures[\rmfamily]{Ligatures=TeX,Scale=1}
\fi
\usepackage{lmodern}
\ifPDFTeX\else
  % xetex/luatex font selection
\fi
% Use upquote if available, for straight quotes in verbatim environments
\IfFileExists{upquote.sty}{\usepackage{upquote}}{}
\IfFileExists{microtype.sty}{% use microtype if available
  \usepackage[]{microtype}
  \UseMicrotypeSet[protrusion]{basicmath} % disable protrusion for tt fonts
}{}
\makeatletter
\@ifundefined{KOMAClassName}{% if non-KOMA class
  \IfFileExists{parskip.sty}{%
    \usepackage{parskip}
  }{% else
    \setlength{\parindent}{0pt}
    \setlength{\parskip}{6pt plus 2pt minus 1pt}}
}{% if KOMA class
  \KOMAoptions{parskip=half}}
\makeatother
\usepackage{longtable,booktabs,array}
\usepackage{caption}
\captionsetup[table]{skip=6pt}
\newcounter{none} % for unnumbered tables
\usepackage{calc} % for calculating minipage widths
% Correct order of tables after \paragraph or \subparagraph
\usepackage{etoolbox}
\makeatletter
\patchcmd\longtable{\par}{\if@noskipsec\mbox{}\fi\par}{}{}
\makeatother
\setlength{\emergencystretch}{3em} % prevent overfull lines
\providecommand{\tightlist}{%
  \setlength{\itemsep}{0pt}\setlength{\parskip}{0pt}}
\usepackage{needspace}
\clubpenalty=10000
\widowpenalty=10000
\pretocmd{\section}{\Needspace{6\baselineskip}}{}{}
\pretocmd{\subsection}{\Needspace{10\baselineskip}}{}{}
\usepackage{bookmark}
\IfFileExists{xurl.sty}{\usepackage{xurl}}{} % add URL line breaks if available
\urlstyle{same}
% fallback for those not using the hyperref driver hyperxmp:
\makeatletter
\@ifundefined{xmpquote}{\newcommand{\xmpquote}[1]{#1}}{}
\makeatother
\hypersetup{
  pdftitle={JAGR-Infinite: Verifying Dependency Resolution in npm Layouts},
  pdfauthor={Bhuvan Prakash · Independent Researcher},
  hidelinks,
  pdfcreator={LaTeX via pandoc}}

\title{JAGR-Infinite: Verifying Dependency Resolution in npm Layouts}
\usepackage{etoolbox}
\makeatletter
\providecommand{\subtitle}[1]{% add subtitle to \maketitle
  \apptocmd{\@title}{\par {\large #1 \par}}{}{}
}
\makeatother
\subtitle{From version bitmaps to witnesses for the package each
consumer can see}
\author{Bhuvan Prakash · Independent Researcher}
\date{13 November 2026 · Jhol Research Series IV}

\begin{document}
\maketitle

\subsection*{Abstract}\label{abstract}
\addcontentsline{toc}{subsection}{Abstract}

A package resolver can choose versions that satisfy the manifest and
still produce a layout in which a consumer sees the wrong dependency.
This paper studies the boundary between selection and placement in
Jhol's JAGR-Infinite resolver. I trace the chart solver, instance
resolver and lockfile writer, and use native Rust probes to distinguish
arithmetic defects from search and representation limits. The baseline
disagreed with 26 of 81 bitmap intersections and missed a valid older
release. A follow-up with corrected bitmap operations and
release-conditioned search agreed with all 81 intersections and
recovered the older release in both selectors. I propose a witness that
records the first visible provider and the empty paths checked before
it. In an exhaustive model of three lookup locations and two versions,
reusing a witness while its checked prefix stayed unchanged produced no
wrong verdicts across 2,187 layout-and-range comparisons. Checking only
the old provider produced 26 false accepts. The model identifies a
concrete way to verify and incrementally recheck a layout. It does not
establish a general npm compatibility guarantee or a performance gain
for the proposed verifier.

\section{Introduction}\label{introduction}

When I look at a Jhol lockfile, I want to answer a question a list of
versions cannot answer: \textbf{which copy will this consumer actually
load?} A project can need two releases of the same package. Both
releases can appear in a manifest inventory while one consumer still
resolves to the wrong copy.

Paper I studies what a resolver may discard from its input. This paper
studies what it must prove about its output. It makes three
contributions:

\begin{enumerate}
\def\labelenumi{\arabic{enumi}.}
\tightlist
\item
  An implementation analysis of where JAGR-Infinite chooses a release,
  where Jhol places it, and where the installer checks it.
\item
  Native probes that expose failures hidden by aggregate install
  results, alongside an independent finite test of first-provider
  witness reuse.
\item
  A verification design that tracks both the provider used by an edge
  and the absence of closer providers. The design specifies which checks
  can be reused after a layout edit.
\end{enumerate}

\section{Related Work}\label{related-work}

Arborist documents the ideal tree and installation layout for npm
{[}1{]}. Its behavior is a useful compatibility baseline under matched
inputs. PacSolve and MaxNPM formulate dependency consistency and
optimization as constraints {[}2{]}. A comparison with them should use
the intersection of their semantic model and Jhol's, rather than
treating different but valid version choices as failures. PubGrub
derives incompatibilities and uses conflict-driven backtracking {[}3{]}.
It provides a useful solver-design comparison. Its default package
identity model needs adaptation before comparing npm placement behavior.
CUDF-based modular package managers separate solver inputs and
objectives from materialization {[}4{]}. Here the verification object is
a concrete path-keyed layout and the provider visible to each consumer.

\section{Implementation Under Study}\label{implementation-under-study}

\subsection{From the command to a
layout}\label{from-the-command-to-a-layout}

Jhol has several ways to satisfy an install request. An eligible
prepared or lockfile layout can bypass new resolution. An unlocked
native install gathers roots and metadata, expands a flat selection, and
can ask JAGR-Infinite's chart for an alternative. The installer checks a
chart candidate against the declared ranges of selected releases before
accepting it. When the flat result cannot represent a conflict, the
installer can attempt a single-pass nested layout from available
metadata. If that route is ineligible or rejected, the instance resolver
is tried. The ordinary path can fall back again. Policy-bound paths
shown in the source fail rather than silently substituting a present-day
solve.

The \href{../../src/main.rs}{CLI dispatch},
\href{../../crates/jhol-core/src/install.rs}{installer} and
\href{../../crates/jhol-core/src/lockfile_write.rs}{lockfile writer}
implement distinct parts of that decision. A successful install alone
does not show which solver handled the graph. A resolver benchmark needs
to record whether its candidate was selected, rejected, bypassed, or
replaced by fallback.

\subsection{Candidate domains and
selection}\label{candidate-domains-and-selection}

The \href{../../crates/jhol-core/src/jagr_infinite.rs}{chart
implementation} maps versions to local ordinals. Candidate masks are in
the target package's ordinal space. Each bit needs its base offset to
have a meaning. The production virtual machine propagates masks and
searches under a budget. In the follow-up chart, each edge carries a
second mask in the parent package's ordinal space, marking the releases
that declare it. The flat selection helper returns either a map keyed by
package name or no candidate. Topological quantities are computed on a
diagnostic branch and do not decide the ordinary virtual-machine
selection.

Let \(D_p\) be the allowed releases of a parent slot \(p\). For a child
slot \(c\), let \(S_{p,c}(v)\) be the child values compatible with
parent release \(v\). A release without that edge imposes no
restriction, so \(S_{p,c}(v)=V_c\), the full child universe for the
model. When several parent slots truly constrain the \textbf{same} child
slot, a sound narrowing step has the form

\[D_c' = D_c \cap \bigcap_p \left(\bigcup_{v\in D_p} S_{p,c}(v)\right).\]

The inner union preserves alternatives. The outer intersection applies
only to a shared child slot. This is a condition for the model, not a
description of every existing propagation step. Once a release is
selected, its exact declarations must be checked. Otherwise, a child
supported by a different candidate parent release can survive. It is
also necessary to represent a candidate release that has no edge at all.

The implementation now uses the mandatory-target part of this rule. It
narrows a target only when the declaring releases cover the parent's
current domain. Search activates a required edge when the chosen parent
release declares it, then checks any already selected child against its
own mask. Backtracking rolls back both active packages and narrowed
domains. The acceptance check recomputes the required closure from the
roots and rejects an incompatible, missing, or unreachable assignment.
It applies to fresh results, continuation candidates and cached
assignments.

\subsection{Instance construction and
placement}\label{instance-construction-and-placement}

The \href{../../crates/jhol-core/src/jagr_nest.rs}{nested resolver}
reads package metadata on demand, orders the worklist deterministically,
and reuses an existing satisfying instance or selects a new one. It
supports multiple major versions. The instance tracks both registry
identity and the name under which it is installed. Its bounds, missing
metadata, and absence of an eligible release have separate outcomes. Its
local preference rule does not systematically backtrack to a different
parent release when that parent's transitive requirements fail.

The \href{../../crates/jhol-core/src/layout.rs}{path layout code} has a
first-visible lookup and a required-edge check. A separate nested
planning path chooses hoisted winners and assigns other instances under
dependents. Its placement map keys by package and version in the
examined source, even though resolver records can distinguish aliases.
Its peer check can look at versions present anywhere in the instance
set. Those are reasons to check alias collisions and peer visibility at
concrete paths before claiming universal placement correctness. The
present paper does not infer an end-to-end failure from those source
risks.

The final lockfile carries paths and artifact provenance. Tarball
integrity and the on-disk layout are separate checks.
\href{../../crates/jhol-treecheck/src/lib.rs}{Treecheck} compares the
installed tree with the authorized layout. It does not hash every
installed JavaScript file. The witness studied here verifies
dependency-edge placement, before byte integrity or lifecycle effects
enter the contract.

\section{First-Provider Witnesses}\label{first-provider-witnesses}

\subsection{Contract}\label{contract}

Fix an immutable metadata snapshot, layout lookup rules, and the
platform and policy inputs. A layout \(L\) maps a canonical slot to
either an identified package instance or absence. For a consumer at path
\(q\) and requested name \(n\), let

\[C(q,n)=(c_1,c_2,\ldots,c_k)\]

be the candidate slots in lookup order. For a required edge \(e\),
validity means that the \textbf{first occupied} slot supplies the
requested package identity and a version satisfying the selected
consumer's declared range. An incompatible nearer provider shadows any
farther compatible provider. For the formal result below, lookup order,
the consumer declaration and the identity/version interpretation are
fixed. Optional, peer, alias and workspace rules require their own
explicit checks before expanding the contract.

\subsection{Proposition and proof}\label{proposition-and-proof}

Let \(j\) be the first occupied slot in \(C(q,n)\), or \(k\) when all
slots are empty. The witness \(W_e\) contains the facts
\(L(c_1),\ldots,L(c_j)\), along with the consumer declaration and policy
context. If those inputs and all recorded slot facts are unchanged in
\(L'\), the edge verdict is unchanged.

\textbf{Proof.} When a provider exists, every recorded predecessor
remains empty and the recorded provider remains the first occupied slot
with the same identity and version. The declaration and rules used to
judge it are unchanged, so the result is unchanged. When no provider
exists, all candidate locations were recorded as empty and remain empty.
This covers successful and failed required lookups. The proof also
explains the need to record absence. A newly occupied nearer location
can change the verdict while leaving the old provider untouched.

A proposed implementation indexes each recorded slot by the edges that
depended on it. For changed slots \(\Delta\) and changed declarations
\(E_{\mathrm{new}}\), recheck

\[A(\Delta)=E_{\mathrm{new}}\;\cup\!
\bigcup_{c\in\Delta}\{e:c\in W_e\}.\]

Here \(c\in W_e\) means the witness recorded a fact about that slot,
including absence. Rule, platform or policy changes invalidate the
relevant context. The reverse index must be updated on every recheck.
Work for those lookups is proportional to the sum of their lookup
lengths, plus change detection and index maintenance. Reading the entire
tree to find changes can dominate the saved checks. This is a design to
implement and evaluate in Jhol.

\section{Evaluation Methodology}\label{evaluation-methodology}

The evaluation separates three evidence sources. The native probe calls
the two Rust resolver APIs on in-memory synthetic package records and
compares bitmap operations with ordinary integer sets. The witness check
independently enumerates a finite placement space. The performance
observations come from older archived loopback experiments that produced
lockfiles without fetching tarballs. None of these experiments runs a
real-project install for this paper.

\section{Results}\label{results}

\subsection{Native Differential
Probes}\label{native-differential-probes}

I built the baseline example in release mode on macOS arm64 and ran it
twice. The two runs produced identical JSON. It uses only synthetic
packuments supplied in memory and does not invoke the CLI or the
materializer.

\begin{longtable}[]{@{}
  >{\raggedright\arraybackslash}p{(\linewidth - 4\tabcolsep) * \real{0.3516}}
  >{\raggedright\arraybackslash}p{(\linewidth - 4\tabcolsep) * \real{0.2967}}
  >{\raggedright\arraybackslash}p{(\linewidth - 4\tabcolsep) * \real{0.3516}}@{}}
\caption{Native differential probes and a multi-major
control.}\tabularnewline
\toprule\noalign{}
\begin{minipage}[b]{\linewidth}\raggedright
Probe
\end{minipage} & \begin{minipage}[b]{\linewidth}\raggedright
Reference expectation
\end{minipage} & \begin{minipage}[b]{\linewidth}\raggedright
Native result
\end{minipage} \\
\midrule\noalign{}
\endfirsthead
\toprule\noalign{}
\begin{minipage}[b]{\linewidth}\raggedright
Probe
\end{minipage} & \begin{minipage}[b]{\linewidth}\raggedright
Reference expectation
\end{minipage} & \begin{minipage}[b]{\linewidth}\raggedright
Native result
\end{minipage} \\
\midrule\noalign{}
\endhead
\bottomrule\noalign{}
\endlastfoot
Same-base bitmap intersection & 81 exact set intersections & 26
disagreed \\
Wide mask intersected with first word & Ordinals 0 to 63, 64 values &
Ordinals 0 to 127, 128 values \\
Conditional dependency & parent 1.0.0 is a valid complete choice & Flat
helper declined. Nested helper chose parent 1.1.0 and reported
Unsatisfiable for child \^{}2.0.0 \\
Two incompatible child majors & Four instances including two child
versions & Flat helper declined. Nested helper returned all four
instances \\
\end{longtable}

\textbf{Bitmap boundary.} The probe takes all ordered pairs of intervals
starting at ordinal zero with upper endpoints 0, 62, 63, 64, 65, 126,
127, 128 and 129. It compares membership with the ordinary intersection.
The equal-base fast path intersects zipped words and leaves the
remaining left-hand words unchanged when the right mask is shorter. The
64-value example shows exactly why that matters. A later validator can
reject some bad selections, but that cannot restore a correct primitive
or prove all callers catch the error. The probe does not show an
incorrect package reaching disk.

\textbf{Conditional dependency.} The root admits parent 1.0.0 and 1.1.0.
Version 1.0.0 has no child requirement. Version 1.1.0 requires child
\^{}2.0.0, but the supplied snapshot has only child 1.0.0. The nested
helper tries 1.1.0, reports a conflict for its child, and does not
revisit the parent. The flat helper also returns no candidate. The
conflict describes the explored branch. It cannot certify the root
problem infeasible when parent 1.0.0 is a witness. The multi-major
control shows a different outcome: the flat representation declines a
layout that the nested instance representation can express.

\subsection{Follow-up Native Probes}\label{follow-up-native-probes}

\begin{longtable}[]{@{}
  >{\raggedright\arraybackslash}p{(\linewidth - 4\tabcolsep) * \real{0.3516}}
  >{\raggedright\arraybackslash}p{(\linewidth - 4\tabcolsep) * \real{0.2967}}
  >{\raggedright\arraybackslash}p{(\linewidth - 4\tabcolsep) * \real{0.3516}}@{}}
\caption{Results for identical synthetic inputs before and after the
resolver changes.}\tabularnewline
\toprule\noalign{}
\begin{minipage}[b]{\linewidth}\raggedright
Probe
\end{minipage} & \begin{minipage}[b]{\linewidth}\raggedright
Baseline
\end{minipage} & \begin{minipage}[b]{\linewidth}\raggedright
Follow-up
\end{minipage} \\
\midrule\noalign{}
\endfirsthead
\toprule\noalign{}
\begin{minipage}[b]{\linewidth}\raggedright
Probe
\end{minipage} & \begin{minipage}[b]{\linewidth}\raggedright
Baseline
\end{minipage} & \begin{minipage}[b]{\linewidth}\raggedright
Follow-up
\end{minipage} \\
\midrule\noalign{}
\endhead
\bottomrule\noalign{}
\endlastfoot
Bitmap intersections & 26 disagreements in 81 pairs & 0 disagreements in
81 pairs \\
Conditional dependency & Neither selector recovered the witness & Both
returned parent 1.0.0 \\
Incompatible child majors & Four nested instances and no flat choice &
Four nested instances and no flat choice \\
\end{longtable}

The release-conditioned flat selector also agreed with exhaustive
feasibility in 81 two-package worlds, including cycles, and each
returned witness satisfied its active edges. The cached-assignment test
changes a release declaration and checks that the previous assignment is
rejected and the cache identity changes. These tests establish results
for the tested flat model. The lookup witness below concerns placement
after selection.
\href{../architecture/experiments/jagr-infinite-native-probe-release-edges-2026-11-13.json}{Follow-up
native results}

\subsection{Native Placement and Frozen
Replay}\label{native-placement-and-frozen-replay}

The baseline CLI exited successfully on a synthetic local-registry
project while one root alias lacked a manifest. I reran the same fixture
after keeping provider instance IDs through path assignment. The
manifest names and versions matched at all six expected paths, including
two root aliases of the same registry package and a nested alias whose
registry target differs from the root alias. After removing the live
tree, a frozen offline install recreated the same six identities with
zero registry requests. The local fixture exercises publication and
lockfile replay. It does not estimate latency or establish compatibility
across arbitrary registries or filesystems.
\href{../../scripts/bench/jagr_nested_install_probe.py}{Probe and
recorded results},
\href{../architecture/experiments/jagr-nested-install-probe-2026-11-13.json}{Raw
observations}

\subsection{Exhaustive Witness Check}\label{exhaustive-witness-check}

Each of three ordered slots can be empty or contain version 1 or 2. The
reference checker enumerates all 27 layouts before and after an edit and
tests ranges accepting version 1, version 2, or either. For each of the
2,187 combinations, it compares a full new lookup with two reuse rules.
The first requires the entire checked prefix to remain unchanged. The
second checks only whether the old provider remains in its old slot.

\begin{longtable}[]{@{}
  >{\raggedright\arraybackslash}p{(\linewidth - 6\tabcolsep) * \real{0.3875}}
  >{\raggedleft\arraybackslash}p{(\linewidth - 6\tabcolsep) * \real{0.2125}}
  >{\raggedleft\arraybackslash}p{(\linewidth - 6\tabcolsep) * \real{0.2250}}
  >{\raggedleft\arraybackslash}p{(\linewidth - 6\tabcolsep) * \real{0.1750}}@{}}
\caption{Exhaustive results for the three-slot, two-version
model.}\tabularnewline
\toprule\noalign{}
\begin{minipage}[b]{\linewidth}\raggedright
Reuse rule
\end{minipage} & \begin{minipage}[b]{\linewidth}\raggedleft
Verdicts reused
\end{minipage} & \begin{minipage}[b]{\linewidth}\raggedleft
Wrong reused verdicts
\end{minipage} & \begin{minipage}[b]{\linewidth}\raggedleft
False accepts
\end{minipage} \\
\midrule\noalign{}
\endfirsthead
\toprule\noalign{}
\begin{minipage}[b]{\linewidth}\raggedright
Reuse rule
\end{minipage} & \begin{minipage}[b]{\linewidth}\raggedleft
Verdicts reused
\end{minipage} & \begin{minipage}[b]{\linewidth}\raggedleft
Wrong reused verdicts
\end{minipage} & \begin{minipage}[b]{\linewidth}\raggedleft
False accepts
\end{minipage} \\
\midrule\noalign{}
\endhead
\bottomrule\noalign{}
\endlastfoot
Checked prefix, including absence & 549 & 0 & 0 \\
Previously present provider only & 702 & 52 & 26 \\
\end{longtable}

The 26 false accepts occur when the old lookup was valid and a change
makes the new lookup invalid while the old provider remains. The other
26 wrong reuses go in the opposite direction. The two reuse rules have
different eligibility, including how they handle an originally absent
provider. The counts measure this finite enumeration, not the rate of
shadowing in npm projects.

\subsection{Archived Performance
Observations}\label{archived-performance-observations}

\begin{longtable}[]{@{}
  >{\raggedright\arraybackslash}p{(\linewidth - 6\tabcolsep) * \real{0.3506}}
  >{\raggedleft\arraybackslash}p{(\linewidth - 6\tabcolsep) * \real{0.1818}}
  >{\raggedleft\arraybackslash}p{(\linewidth - 6\tabcolsep) * \real{0.1558}}
  >{\raggedright\arraybackslash}p{(\linewidth - 6\tabcolsep) * \real{0.3117}}@{}}
\caption{Metadata-only command timings from the archived experiment.
Includes startup and lockfile writing.}\tabularnewline
\toprule\noalign{}
\begin{minipage}[b]{\linewidth}\raggedright
Loopback lockfile workload
\end{minipage} & \begin{minipage}[b]{\linewidth}\raggedleft
Jhol median
\end{minipage} & \begin{minipage}[b]{\linewidth}\raggedleft
Bun median
\end{minipage} & \begin{minipage}[b]{\linewidth}\raggedright
Scope
\end{minipage} \\
\midrule\noalign{}
\endfirsthead
\toprule\noalign{}
\begin{minipage}[b]{\linewidth}\raggedright
Loopback lockfile workload
\end{minipage} & \begin{minipage}[b]{\linewidth}\raggedleft
Jhol median
\end{minipage} & \begin{minipage}[b]{\linewidth}\raggedleft
Bun median
\end{minipage} & \begin{minipage}[b]{\linewidth}\raggedright
Scope
\end{minipage} \\
\midrule\noalign{}
\endhead
\bottomrule\noalign{}
\endlastfoot
1,200 independent roots & 264.7 ms & 84.3 ms & Three runs per tool,
direct selections checked \\
1,200-package chain & 947.7 ms & 311.3 ms & Separate fixture, three runs
per tool \\
\end{longtable}

A separate diagnostic on the chain recorded 498.3 ms fetching and
expanding metadata, 7.1 ms building the chart, and 5.5 ms solving plus
initial materialization. The full resolver path was 531.2 ms in that
run. If fetch and expansion remain unchanged on its critical path,
making all other resolver work free could improve that diagnostic by at
most \(531.2/498.3\approx1.066\). This is a bound for that stage split,
not a speedup measurement. It motivates profiling metadata and
allocation before optimizing solver arithmetic for this workload. The
\href{../research/JAGR_INFINITE_MATH_ONLY_BENCHMARK_RECORD.md}{experiment
record} gives the historical versions, conditions, scheduler trials, and
limitations.

\section{Threats to Validity}\label{threats-to-validity}

The native fixtures are deliberately small. They expose a counterexample
or validate a control, and they cannot measure defect prevalence. The
proposed verifier's proof holds for a fixed snapshot, declarations and
lookup semantics. The executable model uses abstract slots, without
filesystem symlinks, npm peer rules or scripts. The historical
performance record has three repetitions per comparison and uses
generated metadata on loopback. An end-to-end evaluation needs a frozen
corpus, route labels, path-level checks and a second platform.

\section{Conclusion}\label{conclusion}

JAGR-Infinite already gives Jhol a bitmap chart, bounded search, and a
way to resolve multiple releases into instances. The research
opportunity is to make the step from those choices to the layout
independently checkable. The key fact for each required edge is the
first provider the consumer can reach. Its witness must remember what
was absent on the way there. That gives Jhol a concrete condition under
which it can reuse a verdict after an edit.

\section*{Artifact Availability}\label{artifact-availability}
\addcontentsline{toc}{section}{Artifact Availability}

The \href{../../crates/jhol-core/examples/jagr_paper_iv_probe.rs}{native
probe},
\href{../architecture/experiments/jagr-infinite-native-probe-2026-11-13.json}{native
observations},
\href{../architecture/experiments/jagr-infinite-native-probe-release-edges-2026-11-13.json}{follow-up
native observations}, and
\href{../../scripts/bench/jagr_iv_witness.py}{finite witness checker}
accompany this manuscript. The Rust probe runs with Cargo's locked,
offline release build of the jhol-core example, with the http-reqwest
feature and selection caching disabled. The witness checker requires
Python 3 and no packages. See the \href{../../LICENSE}{repository
license} for the software terms.

\Needspace{13\baselineskip}
\section*{References}\label{references}
\addcontentsline{toc}{section}{References}

\begin{enumerate}
\def\labelenumi{\arabic{enumi}.}
\tightlist
\item
  npm CLI. ``Arborist: Build Ideal Tree.''
  \href{https://github.com/npm/cli/blob/latest/workspaces/arborist/docs/ideal-tree.md}{Design
  documentation}.
\item
  D. Pinckney, F. Cassano, A. Guha, J. Bell, M. Culpo, and T. Gamblin.
  ``Flexible and Optimal Dependency Management via Max-SMT.''
  \emph{IEEE/ACM International Conference on Software Engineering},
  2023. \href{https://arxiv.org/abs/2203.13737}{Paper and preprint}.
\item
  Dart pub team. ``PubGrub Version Solving Algorithm.''
  \href{https://dart.googlesource.com/pub.git/+/f27dcfdb/doc/solver.md}{Design
  documentation}.
\item
  P. Abate, R. Di Cosmo, R. Treinen, and S. Zacchiroli. ``A Modular
  Package Manager Architecture.'' \emph{Information and Software
  Technology}, 2013.
  \href{https://www.dicosmo.org/Articles/2013-AbateDiCosmoTreinenZacchiroli-Ist.pdf}{Paper}.
\end{enumerate}


\section*{Post-publication corrections (measured)}

\textbf{The resolver is measured against npm, and the first measurement disagreed.} A ten-scenario
compatibility suite (trees compared as \code{name@version} sets) initially found divergence on four:
\pkg{express}, \pkg{jest}, \pkg{sharp} and \pkg{@babel}. Root causes, each fixed at source:
(i) a build-time snapshot answering exact pins while online; (ii) version lists pruned by
constraints-seen-so-far; (iii) a failing chart solve falling back silently; (iv) versions changed
after the dependency walk, leaving dependencies unfetched; (v) selected peers pruned together with
their subtrees; (vi) a wildcard test applied to the whole version string, so any prerelease
containing the letter \code{x} (\code{-no-external-plugins}) was read as a range.

\textbf{Inputs must be authoritative or pinned.} Resolution now reads the full packument; the
abbreviated document is retained for prefetching, where a stale hint cannot decide a version. Both
cache layers share one freshness policy. \code{--as-of} keeps publish dates, without which a dated
resolution silently produces today's tree.

\textbf{Correction: absence of a document is not a verdict about a package.} An optional
dependency whose packument had not been fetched was silently dropped (\code{if optional \{ continue \}}),
which lost \pkg{fsevents} from a \pkg{jest} tree on macOS while npm installed it. A missing document
is now reported as a fetch request; an optional edge is dropped only after the registry has been
asked and does not have the package.

\textbf{Status.} All ten rows match npm on macOS arm64; Linux and Windows are unexecuted, and
lifecycle scripts are not yet part of the comparison.


\section{Post-publication corrections: peer edges and platform identity (measured)}

Three defects shared one cause: \emph{a tree fact inferred from the wrong evidence}. Each is recorded
with the measurement that found it, because in every case the tree looked correct while its own plan
disagreed.

\begin{enumerate}
\item \textbf{Peer edges were advisory, not structural.} Installing \texttt{@testing-library/react\@14}
left \texttt{@types/react-dom} in the tree without its peer \texttt{@types/react}. The resolver
detected the requirement and printed ``got absent'' --- so the graph knew and the layout did not. A
post-pass over resolved instances cannot repair this: an instance that no edge refers to never
reaches the layout. The correction is that peers are \emph{demands}, queued at the same point as
dependencies so the ordinary edge bookkeeping places them.
\item \textbf{Optional peers are not demands.} Queueing every peer made \texttt{ts-node}'s optional
peers \texttt{@swc/core} and \texttt{@swc/wasm} real packages. The lockfile then listed
\texttt{@swc/wasm}, the layout refused it as a WebAssembly build, and \texttt{verify} failed with
\texttt{node_modules/@swc/wasm is missing} --- the tree contradicting its own plan. Optionality is
\texttt{peerDependenciesMeta}, and it must be read before the demand is queued, not after.
\item \textbf{Platform identity is declared metadata, never a name.} The materialisation filter
rejected packages by tokenising their names, and \texttt{@swc/wasm} contains the token \texttt{wasm}.
The package is real, is installed on every platform, and was silently dropped from trees that listed
it. The asymmetry settles the rule: skipping a loadable package breaks a tree, installing an
unloadable one wastes bytes. Only \texttt{os}/\texttt{cpu}/\texttt{libc} fields may exclude an entry,
which is why npm writes them. Name-token inference survives only where no fields exist at all.
\end{enumerate}

One measurement worth keeping: the first run after the metadata fix still dropped the entry, because
parsed lockfile entries are cached in a binary keyed by lockfile content and platform \emph{only}. A
cache keyed by its input alone will happily return answers produced by an older decision procedure,
so the key must carry the rule version.


\section{Post-publication correction: a lockfile states intent, not just an answer}

Measured on a real tree, the writer emitted \textbf{55 empty \texttt{dependencies} maps out of
70 entries} and pinned the rest to wired versions. Two rules follow, and both are about what a
lockfile \emph{is}.

First, every field that describes a package is \emph{declared} metadata: dependencies, optional
dependencies, peers, peer optionality, platform. It is copied from the published version document
by one constructor, never assembled by hand, because a hand-built entry at forty sites is forty
chances to write an empty map.

Second, a writer may consult a cache but must never derive its content from one. The lookup was
TTL-gated at sixty seconds, so two entries of the same file disagreed about the same question:
\texttt{accepts} came out correct and \texttt{express} came out empty because its packument had
aged out. A published version is immutable --- only \texttt{dist-tags} and \texttt{time} move ---
so the immutable half is read without a freshness window, and the lookup is a \emph{fallback}:
a miss must not erase edges the tree already knows.

Result: 68 of 69 entries match npm's dependency maps byte for byte. The exception is an optional
dependency recorded in \texttt{dependencies} rather than \texttt{optionalDependencies}, because
the entry type has no optional-edge class; that is a stated gap, not an unmeasured one.


\section{Post-publication correction: what a lockfile entry must contain}

After `dependencies` became declared ranges, the entry \emph{shape} still disagreed with npm. The
measurement was done field for field on a real tree, and it took three rounds:

\begin{center}
\begin{tabular}{lcl}
identical entries & change & why \\
\hline
0 / 69 & -- & empty \texttt{dependencies} maps, no \texttt{license} \\
48 / 69 & facts carried on the entry & a writer must not derive facts from a cache \\
68 / 69 & \texttt{license} retained by pruning & a version fact that the prune removed \\
69 / 69 & install-script checks all three hooks & \texttt{esbuild} works in \texttt{postinstall} \\
\end{tabular}
\end{center}

Two principles follow. First, \textbf{facts travel with the entry}: every declared field is copied
from the published version document once, at construction, and read from the entry afterwards.
Second, \textbf{optionality is a class, not a merge}: optional edges belong under
\texttt{optionalDependencies} and must not be copied into the required set, which is precisely what
made the last entry differ.

\end{document}
